Na camada de segurança, três frentes ficam abertas. A primeira é a avaliação em uma rede LoRaWAN completa, com \textit{Network Server} e \textit{Application Server}, que é o cenário em que a lacuna de integridade fim a fim discutida no Capítulo~\ref{cap:especificacao} de fato se manifesta, e em que o custo adicional da camada, os 20 bytes de contador e \textit{tag} por mensagem, pode ser medido em tempo no ar e energia de transmissão em vez de calculado. A segunda é a migração para o Ascon-AEAD128 do SP 800-232 quando houver implementação estável para o ESP32, com a confirmação por medição de que o custo não muda. A terceira é a avaliação de canais laterais além do tempo, por consumo de potência e emissão eletromagnética. Vale ainda estudar o provisionamento de chaves em escala, já que o procedimento adotado, uma chave gerada em cada nó e copiada ao receptor na bancada, atende a um protótipo com poucos nós e não a uma implantação com dezenas deles.

%%%%%%%%%%%%%%%%%%%%%%%%%%%%%%%%%%%%%%%%%%%%%%%%%%%%

\subsection{Segurança Física do ESP32}
\label{subsec:seguranca_fisica}
A implementação Ascon-128a fornece confidencialidade e integridade para os dados em trânsito, mas a integridade do sistema como um todo depende do sigilo da chave criptográfica armazenada no nó sensor. Como dispositivos de borda (\textit{Edge}) operam em ambientes fisicamente hostis e não supervisionados, são alvos potenciais para ataques de extração de memória (\textit{Flash Dumping}).

O escopo desta pesquisa concentrou-se na eficiência das primitivas criptográficas. Para uma implantação real, porém, a segurança em nível de \textit{hardware} é uma camada indispensável.

Para mitigar essas vulnerabilidades físicas em desenvolvimentos futuros, recomenda-se a adoção dos mecanismos nativos de segurança da família ESP32 \cite{espressif2024flash,kosemen2021tamper}. A chave simétrica do Ascon pode ser protegida pela criptografia da partição não volátil (NVS), cujas chaves ficam em uma partição própria cifrada pela \textit{Flash Encryption}, com a chave mestre gravada nos \textit{eFuses}, de programação única (OTP, \textit{One-Time Programmable}), e protegida contra leitura externa (\textit{Read Protect}). A Inicialização Segura (\textit{Secure Boot}), que verifica a assinatura do \textit{firmware} a cada inicialização, previne regravação maliciosa. Com esses mecanismos, a eficiência criptográfica do Ascon passaria a operar sobre uma base de \textit{hardware} que protege os dados em repouso e o código contra acesso físico não autorizado.
